% !PRINTED 611
\weberpage{611}
%<<CONT
auszuschliessen ist, weil sonst ja durch (3) alle verschiedenen $z$ in dasselbe $b$
übergehen, also gar keine Permutation der Grössen $x_z$ ausgedrückt wäre. Ist aber $a$
von 0 verschieden $(\mathrm{mod}\ n)$, so wird durch (3) immer eine Permutation
dargestellt sein, weil dann nur, wenn $z \equiv z'$ ist, $az+b \equiv az'+b$ sein kann.
Wir nennen (3) eine Substitution für $z$, die unter den $x$ eine Permutation
hervorbringt. Die Anzahl der verschiedenen linearen Substitutionen von der Form (3)
ist $n(n-1)$. Ihnen entsprechen ebenso viele verschiedene Permutationen unter den $x$;
denn es kann nur dann für jedes $z$
\[ az + b \equiv a'z + b' \, (\mathrm{mod}\ n) \]
sein, wenn $b \equiv b'$ (aus $z = 0$ zu schliessen) und $a \equiv a'$ (aus $z = 1$ zu
schliessen).

Die Gesammtheit der Permutationen, die durch (3) dargestellt sind, bildet eine Gruppe;
denn es seien
\[ \lambda = (z,\ az+b), \qquad \lambda' = (z,\ a'z+b') \]
zwei von diesen Permutationen, so ist
\[
\lambda \lambda' =
\begin{pmatrix} x_z \\ x_{az+b} \end{pmatrix}
\begin{pmatrix} x_z \\ x_{a'z+b'} \end{pmatrix}
=
\begin{pmatrix} x_z \\ x_{a'(az+b)+b'} \end{pmatrix},
\]
also
\weq{4}{\lambda \lambda' = \lambda'' = (z,\ aa'z + a'b + b'),}
was wieder von der Form $(z, a''z+b'')$ ist.

Die durch alle Substitutionen von der Form (3) gebildete Gruppe, die eine
Verallgemeinerung der cyklischen Gruppe (2) ist, wollen wir eine
\gesperrt{lineare Gruppe} nennen\footnote{\gesperrt{Kronecker} nennt nur diese Gruppe
metacyklisch.}; die in ihr enthaltenen Permutationen sollen
lineare Permutationen heissen.

In der linearen Gruppe sind verschiedene Divisoren enthalten. Wir finden darunter eine
Gruppe vom Grade $n-1$, nämlich $(z, az)$, die selbst wieder Divisoren haben kann.

Die Compositionsregel (4) zeigt, wie wir alle Divisoren der linearen Gruppe bilden
können. Ist $\lambda = (z, az+b)$ eine lineare Substitution, so ergiebt sich nach (4)
durch Wiederholung für jeden Exponenten $h$
\[ \lambda^h = [z,\ a^h z + (1 + a + \cdots + a^{h-1})b], \]
und für $a = 1$
\[ \lambda^h = (z,\ z + hb). \]
